\chapter*{Záver}  % chapter* je necislovana kapitola
\addcontentsline{toc}{chapter}{Záver} % rucne pridanie do obsahu
\markboth{Záver}{Záver} % vyriesenie hlaviciek

Pôvodný zámer bakalárskej práce bol vytvoriť podporný systém pre manažérov KIB, ktorý nemajú mnoho znalostí v oblasti KIB.
Tento systém im mal pomôcť naplniť legislatívne požiadavky na ochranu informačných systémov ISVS pod ich
správou\footnote{Ciele práce sme rozoberali v Úvode a aj v kapitole \ref{ch:after-the-battle}.}.

Počas prípravy bakalárskej práce sa ukázalo, že povinnosti vyplývajúce z legislatívnych noriem(\cite{69/2018},
    \cite{362/2018}, \cite{95/2019} a \cite{179/2020}) bolo podstatne viac, než sme na začiatku predpokladali a vnútorne 
tvorili zložitý a nekonzistentný systém. Ťažisko práce pozostávalo v analytickej činnosti, kde sme potrebovali
identifikovať legislatívne požiadavky a následne ich konkretizovať pomocou technických noriem (\cite{27001} a \cite{27002})
a metodik (\cite{BSI1}, \cite{BSI2}, \cite{BSI3}, \cite{GSK}).

Právne normy obsahovali veľké množstvo všeobecných a technických požiadaviek. Zo všetkých požiadaviek pre správcu
malého ISVS je najdôležitejšie mať manažéra KIB a zaviesť ISMS pokrývajúci všetky informačné systémy organizácie. Tým
dostalo riešenie KIB pevný rámec. Ale aj postup podľa jednoduchého návodu pri zavádzaní ISMS si vyžaduje aspoň základné
znalosti KIB od vedenia organizácie, manažéra KIB a zamestnancov. Bez základných znalostí by nebolo možné ani vytvoriť
odborné dokumenty KIB, ani by sa nedalo očakávať, že tým dokumentom niekto porozumie a bude sa nimi vedieť riadiť.

Vytvorili sme jednoduchý systém, ktorý manažéra KIB oboznámi a vysvetlí mu problematiku KIB, získa od neho informácie o
aktuálnom stave KIB v organizácií a povedie ho pri zavádzaní ISMS\footnote{Dokonca počas tohoto procesu od neho 
získame dôležité informácie pre spísanie rôznych dokumentov KIB.}. Tento \enquote{bakalársky} systém je použiteľný ako
základ profesionálneho systému pre podporu manažérov KIB. Po jeho čiastočnej implementácii sme tiež získali predstavu
o tom, ako by bolo potrebné systém upraviť a rozšíriť, aby vyhovoval potrebám manažéra KIB.

Vytvorenie profesionálneho podporného systému si vyžaduje nie len prax a podrobnejšiu znalosť prostredia, v ktorom sa má
systém použivať, ale aj podstatne viac času, vedomostí a práce, ktorej rozsah a odborná úroveň presahuje rámec
bakalárskej práce a úsilie jedného človeka.



